\chapter{Open Problems and a Research Roadmap}\label{ch:open}

Every earlier chapter answered a question the project had already closed: a theorem stated,
proved, and checked by the kernel. This chapter inverts the exercise. It asks, for roughly
two dozen concrete pieces of mathematics and physics that this book has touched but not
finished, exactly what would have to be stated and proved for the badge \TA\ to apply, what
already-checked result it would build on, how hard that looks, and what is or is not
already in Mathlib. A student should care because the gap between ``this book explains the
idea'' and ``this book proves the theorem'' is precisely where new work happens, and because
a roadmap entry that names its prerequisites and its Mathlib gap is a project a reader can
actually start, not a wish. The chapter ends, as every other has, with a worked example --
here, two of them, one purely a formalization target and one a numeric check a reader can
reproduce by hand -- because a roadmap item is only as good as the first concrete step
someone can take on it.

%=====================================================================================
\section{How to read an entry}\label{sec:open:how}

Every item below carries four pieces of information: the chapter(s) it extends, a
difficulty ($\star$ a term-project exercise, $\star\star$ a paper-sized piece of work,
$\star\star\star$ open-ended), the state of the relevant Mathlib machinery, and a
\emph{precise suggested Lean statement} where one exists. Suggested statements are written
as \texttt{leancode} blocks and are, by the convention this book has used throughout and that
\texttt{tools/check\_book\_lean\_names.py} enforces, marked \textbf{not yet proved}: no such
declaration exists in the compiled environment, and the name is invented for this chapter
only. Every other identifier printed in this book, including in this chapter's prose, names
a real, compiled declaration, quoted from the file that defines it.

\begin{pitfallbox}[A precise false statement is worse than a vague true one]
It is tempting to round a hard theorem down to something easy and call the easy version
``the open problem.'' \cref{sec:lattices:classification}'s classification theorem is Tier
L precisely because Serre's proof uses machinery (the classification of indefinite
quadratic forms via genus theory) this project has not formalized; a roadmap entry that
said ``formalize: every even unimodular lattice of signature $(3,19)$ is isometric to
$3U\oplus2E_8(-1)$'' without naming that gap would be an honest goal but a misleading
project description. Every item below names the specific Mathlib or project machinery it
is missing, not just the headline claim.
\end{pitfallbox}

The items are grouped in four sections: extensions of the kernel-checked lattice and
duality layer (\cref{sec:open:lattice}), targets in double field theory and moonshine
(\cref{sec:open:dft}), engineering debt already located by the atlas
(\cref{sec:open:atlas}), and physics-facing questions the dual-scale proposal must answer
to become more than Tier C (\cref{sec:open:physics}).

%=====================================================================================
\section{Extending the kernel-checked lattice and duality layer}\label{sec:open:lattice}

\begin{center}\small
\begin{tabular}{@{}p{0.6cm}p{5.4cm}p{1.7cm}cp{4.6cm}@{}}\toprule
\# & Project & Prereq. & Diff. & Mathlib / project state \\ \midrule
1 & Narain lattice $\Gamma^{d,d}$ even self-duality for general $d$ & \cref{ch:lattices,ch:narain} & $\star\star$ & \lean{IsEvenDiag}, \lean{IsUnimodular}, \lean{Matrix.fromBlocks} all exist; induction on $d$ via direct sum is new \\
2 & $\Odd{2}$ is generated by factorized dualities, reflections and $\SLZ$ & \cref{ch:odd,ch:torus2} & $\star\star$ & GPR's own $d=2$ statement is not fully proved even in the literature source (\cref{sec:open:duality-group}) \\
3 & The full $O(\Gamma^{4,20})$ action on the Mukai lattice & \cref{ch:mukai} & $\star\star\star$ & only $(-2)$-reflections at simple roots are Tier A (\cref{ex:open:e8weyl}); no group-generation machinery in the project \\
4 & Cartan--Dieudonné: every isometry of $E_8(-1)$ is a product of the 8 simple reflections & \cref{ex:open:e8weyl} & $\star\star\star$ & needs a notion of ``group generated by a finite set of matrices,'' absent from the project; Mathlib has \lean{Subgroup.closure} but no bridge to this setting \\
5 & Buscher rules as a theorem about \texttt{DFT.GeneralizedMetric} for general $d$, not just $d=1$ & \cref{ch:buscher,ch:genmetric} & $\star\star$ & the $d=1$ statement (\cref{ch:buscher}) is an explicit $2\times2$ computation; the general case needs a genuine block decomposition argument \\
\bottomrule\end{tabular}
\end{center}

\subsection{Worked example: what a $d=2$ Narain statement would say}\label{sec:open:narain2}

Item 1 is the most tractable of the five, and working it out by hand shows both why the
step from $d=1$ to $d=2$ is easy and why the step to general $d$ is not. Recall
\lean{hyperbolicU} of \cref{ch:lattices}, the Gram matrix $U=\begin{psmallmatrix}0&1\\1&0
\end{psmallmatrix}$, already proved even (\lean{hyperbolicU\_evenDiag}) and unimodular
(\lean{hyperbolicU\_unimodular}, i.e.\ $\det U=\pm1$; here $\det U=-1$). For $d=2$, take
the block-diagonal Gram matrix
\begin{equation}\label{eq:open:blockU}
  U\oplus U=\begin{pmatrix}0&1&0&0\\1&0&0&0\\0&0&0&1\\0&0&1&0\end{pmatrix}.
\end{equation}
Its diagonal entries are the diagonal entries of the two copies of $U$, all zero, hence
even: \lean{IsEvenDiag} is closed under direct sum for free, by definition, since
$\text{diag}(A\oplus B)$ is just the concatenation of $\text{diag}(A)$ and
$\text{diag}(B)$. Unimodularity is one line of linear algebra:
$\det(A\oplus B)=\det(A)\det(B)$ for block-diagonal matrices, a standard Mathlib fact
about \lean{Matrix.fromBlocks} (the zero off-diagonal specialization of its determinant
lemma), so $\det(U\oplus U)=(-1)(-1)=1$, and \lean{IsUnimodular} follows. The suggested
statement is therefore short:
\begin{leancode}
-- not yet proved
theorem narain_evenDiag_two :
    IsEvenDiag (Matrix.fromBlocks hyperbolicU 0 0 hyperbolicU) := by
  intro i; fin_cases i <;> simp [hyperbolicU, Matrix.fromBlocks]
theorem narain_unimodular_two :
    IsUnimodular (Matrix.fromBlocks hyperbolicU 0 0 hyperbolicU) := by
  simp [IsUnimodular, Matrix.det_fromBlocks_zero21, hyperbolicU_unimodular]
\end{leancode}
Why does this not simply generalize to every $d$ by induction? Because \lean{IsEvenDiag}
and \lean{IsUnimodular} are exactly the two properties that are preserved by direct sum
--- diagonal entries do not interact across the block boundary, and determinant is
multiplicative on block-diagonal matrices --- so the induction step \emph{for these two
properties alone} is genuinely easy, and a full statement for general $d$ (say, packaging
$d$ copies of $U$ as a function \lean{Fin d → Fin 2 → Fin 2 → ℤ} block matrix and proving
both properties by induction on $d$) is item 1 above, difficulty $\star\star$ only because
of the bookkeeping of indexing $d$ copies, not because of new mathematics. What the direct
sum \emph{cannot} produce is a Narain lattice with off-diagonal ``instanton'' shifts, i.e.\
a nontrivial $B$-field background; that case needs the full $\Odd{d}$ machinery of
\cref{ch:odd} and is not covered by this item.

\begin{leanbox}[The prerequisite this item builds on]
\begin{leancode}
theorem hyperbolicU_evenDiag : IsEvenDiag hyperbolicU := by ...
theorem hyperbolicU_unimodular : IsUnimodular hyperbolicU := ...
\end{leancode}
\textbf{Reading.} The hyperbolic plane's Gram matrix has even diagonal (both entries
zero) and determinant $\pm1$ (here $-1$). \textbf{Proof idea.} Both are two-line
computations on the explicit $2\times2$ matrix \source{Hyperbolic.lean, as cited at
\cref{ch:lattices}}. \textbf{Not proved.} Anything about $d>1$ without an explicit
construction like \eqref{eq:open:blockU}; that is exactly item 1. \TA
\end{leanbox}

\subsection{The $d=2$ duality group: what GPR itself does and does not claim}
\label{sec:open:duality-group}

Item 2 looks, at first, like it should already be finished, because \cref{ch:torus2}
proves a $d=2$ result about $\SLZ$. Reading the source pinned there precisely is the whole
point of the exercise. GPR's own statement of the $d=2$ duality group is
\emph{``The duality group is $\SLZ\times\SLZ\otimes_S[\Z_2\times\Z_2]$''}
\source{GPR, "The $d=2$ Example," ll.~7389 ff.}, and this book's own use of that source
(\texttt{docs/PAPER\_REVISION\_BRIEF\_2026\_09\_17.md §3.2}) is explicit that only one
piece of it is Tier A here: \emph{``a basis change with $\det A=1$ commutes with every
$\rho$-translation $g_{tJ}$ (Tier A, this direction only\ldots)\; GPR's `The $d=2$ Example'
is not claimed in full; only $\SLZ_\tau$ commuting with the $\rho$-translations is Tier
A.''} The gap between the one commuting subgroup already proved and the full semidirect
product with $\Z_2\times\Z_2$ GPR states is item 2's actual content: it is not a gap in
this project's diligence but a gap the source itself leaves unproved at the point it is
read, since GPR states the full group as a fact of the CFT without a self-contained
algebraic derivation from the generators used in \cref{ch:odd}. A project attacking item 2
should first locate, inside GPR or a later reference, an actual proof of the semidirect
product structure (as opposed to a statement of it), before attempting a Lean
formalization -- the literature gap must be closed, or worked around by choosing a
narrower, provable sub-statement, before the Lean gap can be.

%=====================================================================================
\section{Double field theory, modular forms and moonshine}\label{sec:open:dft}

\begin{center}\small
\begin{tabular}{@{}p{0.6cm}p{5.4cm}p{1.7cm}cp{4.6cm}@{}}\toprule
\# & Project & Prereq. & Diff. & Mathlib / project state \\ \midrule
6\label{ex:open:dualscale-unique} & Uniqueness of the dual-scale minimizer up to $O(d)$-conjugation & \cref{ch:dualscale} & $\star\star$ & \lean{dualScale\_ge}, \lean{dualScale\_one} exist; the equality case of the trace-of-a-PSD-sandwich argument is not extracted \\
7 & Siegel--Narain theta function, modular transformation law & \cref{ch:narain} & $\star\star\star$ & Mathlib's modular-forms library covers weight-$k$ forms for $\SLZ$, not vector-valued or Siegel forms; a real gap \\
8 & $A_n$ (EOT coefficients) positivity, extended past $n=2$ & \cref{ch:moonshine} & $\star$ & the existing decomposition theorems cover $A_1=45$, $A_2=231$; extending to $A_3,A_4,\dots$ needs the $M_{24}$ character table entries this project has not encoded \\
9 & $A_n$ as a virtual $M_{24}$ character (i.e.\ $A_n=\sum_i m_i\chi_i$ with $m_i\in\Z$) & item 8 & $\star\star\star$ & needs the full 26-dimensional $M_{24}$ character table; \lean{M24RepDim} covers only irrep dimensions, not characters \\
10 & Jacobi form weight/index transformation law for the K3 elliptic genus & \cref{ch:ellgenus} & $\star\star\star$ & no Jacobi-form machinery in Mathlib or the project; \cref{ch:ellgenus} states the genus as a $q,y$-expansion, not as a section of a line bundle \\
11 & Modular invariance of the full Narain torus partition function & \cref{ch:narain,ch:tduality} & $\star\star\star$ & needs item 7; the level-matching and mass-formula pieces (\cref{ch:circle,ch:odd}) exist separately but are never assembled into a partition function \\
\bottomrule\end{tabular}
\end{center}

\subsection{Worked example: the state of $A_n$ positivity}\label{sec:open:moonshine-worked}

Item 8 is the cheapest entry in this chapter to make concrete, because the numbers are
already in this book. The Eguchi--Ooguri--Tachikawa coefficients are tabulated as
\begin{equation}\label{eq:open:antable}
  A_1,A_2,A_3,\dots=45,\,231,\,770,\,2277,\,5796,\dots
\end{equation}
and the paper that first isolated them states plainly, of exactly this sequence,
\emph{``and it was conjectured that they are all positive integers''} \source{EOT,
ll.~208--210}. By 2026 the positivity and integrality of every $A_n$ is a theorem of the
literature (Gannon and others have since given closed forms), so item 8 is not a physics
conjecture but a bookkeeping gap: this project's Lean statements
(\cref{ch:m24,ch:moonshine}) decompose $A_1=45$ and $A_2=231$ into $M_{24}$ irreducible
representation dimensions and check the arithmetic by \lean{decide}, but stop there. A
$\star$ project is to state and prove, in the same style,
\begin{leancode}
-- not yet proved
theorem A3_positive : (770 : ℤ) > 0 := by decide
theorem A3_decomposition :
    770 = 1 + 23 + 253 + 1771 - 1 - 45 - 231 - 770 + 770 := by decide  -- illustrative only
\end{leancode}
where the actual right-hand side must be the correct signed sum of $M_{24}$ irrep
dimensions for $A_3$, which requires consulting the character table (Gaberdiel--Hohenegger
--Volpato, \source{GHV §3.1, ll.~584--700}, already a source key of this book) rather than
guessing coefficients -- exactly the discipline \cref{sec:pipeline:sympy} asks for before
any proof search begins. The honest difficulty rating $\star$ reflects that the arithmetic
step is trivial (\lean{decide} on a finite sum) once the correct decomposition is looked
up; the work is entirely in reading the source correctly, which is why item 9 (deriving
the decomposition rule in general, rather than checking it term by term) is $\star\star\star$.

%=====================================================================================
\section{From unification candidate to bridge: the remaining atlas debt}\label{sec:open:atlas}

\Cref{ch:atlas} refactored three clusters of duplicated theorems in detail ($\chi(K3)=24$
proved six times, the Mukai rank three times, RR tadpole cancellation four times with two
accountings) and described the mechanical repair: pick one canonical definition, replace
every duplicate with a corollary derived from it. The atlas's unification-candidate list
has entries that chapter did not work through; each is a self-contained $\star$ or
$\star\star$ project of exactly the same shape.

\begin{center}\small
\begin{tabular}{@{}p{0.6cm}p{8.9cm}p{2.9cm}c@{}}\toprule
\# & Cluster (cosine similarity, dependency Jaccard) & Libraries & Diff. \\ \midrule
12 & \lean{golay\_self\_dual\_dimension} $\sim$ \lean{golay\_length\_matches\_k3\_euler} $\sim$ \lean{golay\_code\_rate\_is\_half} (0.55, 0.01; 0.54, 0.03) & Lean5Corpus,\newline StringTheory\-\newline Foundation & $\star$ \\
13 & \lean{dissipation\_strictly\_positive} $\sim$ \lean{energy\_dissipation\_monotonic} (0.36, 0.03) & Lean5Corpus,\newline StringTheory\-\newline Foundation & $\star$ \\
14 & \lean{tachyon\_condensation\_endpoint} $\sim$ \lean{tachyon\_condensation\_preserves\_rr\_charge} (0.35, 0.03) & DualScale\-\newline Validation,\newline DualScaleM24\-\newline Formalization & $\star$ \\
15 & \lean{sdc\_mass\_suppression} $\sim$ \lean{sdc\_tower\_suppression} (0.34, 0.01) & DualScale\-\newline Validation,\newline StringTheory\-\newline Formalization & $\star$ \\
16 & \lean{buscher\_log\_involution} (three copies) $\sim$ \lean{buscher\_involution} (0.36, 0.00) & DoubleField\-\newline Theory,\newline DualScale\-\newline Validation,\newline DualScaleM24\-\newline Formalization & $\star\star$ \\
\bottomrule\end{tabular}
\end{center}
\source{atlas.md, ll.~133--134 (item 12), l.~181 (item 13), l.~187 (item 14), l.~189
(item 15), l.~177 (item 16)}

Item 16 is rated $\star\star$ rather than $\star$ because, unlike the others, its
low dependency-Jaccard score together with a fairly high cosine similarity means the four
theorems plausibly encode \emph{different} constructions that happen to share a name and a
one-line statement shape (an involution property) -- the refactor here must start by
reading all four \lean{.lean} files and confirming whether they are the same fact before
choosing a canonical definition, exactly the caution \cref{ch:atlas} applied to the tadpole
cluster, which turned out to hide two different countings behind one number. A reader
attempting this item should treat ``are these the same theorem'' as the first question,
not an assumption.

%=====================================================================================
\section{A faithful, non-arithmetic-shadow formalization}\label{sec:open:shadow}

\cref{sec:lattices:classification}'s honesty rule applies to more than the classification
theorem: several of this project's oldest libraries, flagged throughout this book as
\emph{arithmetic shadows} -- \texttt{DoubleFieldTheory}, \texttt{StringTheoryFoundation},
\texttt{DualScaleM24\-Formalization} and two more -- model a physical quantity by an
integer or rational and prove an arithmetic consequence of that model, not a statement
about the physical object itself. The RR tadpole identity is the clearest case:

\begin{center}\small
\begin{tabular}{@{}p{0.6cm}p{5.4cm}p{1.7cm}cp{4.6cm}@{}}\toprule
\# & Project & Prereq. & Diff. & Mathlib / project state \\ \midrule
17 & RR tadpole cancellation as a statement about a cohomology class $[G]\in H^4(K3\times T^2,\Z)$, not an integer sum & \cref{ch:tadpole,ch:atlas} & $\star\star\star$ & no cohomology in the project or in the relevant part of Mathlib; \texttt{Flux.Integrality} only proves evenness of a Kronecker product \\
18 & The Dirac quantization condition $[G/2\pi]\in H^4(X,\Z)$ as a hypothesis, not an assumed integrality & item 17 & $\star\star\star$ & needs a formal notion of de~Rham/singular cohomology class and its integral lattice; a project-defining gap, not a one-file fix \\
\bottomrule\end{tabular}
\end{center}

\begin{historybox}[Why the shadow was chosen in the first place]
\cref{sec:pipeline:judge} explains the pipeline's operating constraint: goals must be
closable by a 7B local prover or a cheap hosted model before escalating, and Mathlib's
coverage of algebraic topology (singular cohomology of a specific space, cup products,
Poincar\'e duality on a compact 4-manifold) is thin enough that a faithful statement of
items 17--18 was, at the time this book was written, outside what any tier of the pipeline
could attempt. The integer bookkeeping of \texttt{Flux.Tadpole} is not a mistake; it is the
largest true statement the project could currently prove about the physical content
Dasgupta--Rajesh--Sethi state in words \source{DRS §3.1, l.~640; ll.~1265--1275,
1366--1370, already a source of \cref{ch:tadpole}}.
\end{historybox}

%=====================================================================================
\section{What the dual-scale proposal needs to become testable}\label{sec:open:physics}

The remaining items are physics questions, Tier C by the convention fixed for this whole
book (a reading of a Lean object as a physical statement is never Tier A), for which a
Lean statement is not yet the bottleneck -- a precise mathematical \emph{formulation} is.

\begin{center}\small
\begin{tabular}{@{}p{0.6cm}p{7.9cm}p{2.9cm}@{}}\toprule
\# & Question & Nearest chapter \\ \midrule
19 & What observable would distinguish $R+\ap/R\ge2\sqrt{\ap}$ from ordinary sub-Planckian physics? & \cref{ch:cosmology} \\
20 & A precise, checkable invariant matching heterotic on $K3\times T^2$ to type~II on a specific Calabi--Yau threefold & \cref{ch:k3t2} \\
21 & The duality group $\Gamma=\Gamma_W\rtimes_S\Gamma_M$ of a Calabi--Yau moduli space, beyond the one-modulus case & \cref{ch:torus2} \\
22 & An exact count of the $M_{24}$-orbit structure of $24$ points, for symmetry groups beyond the Fermat quartic example & \cref{ch:m24} \\
\bottomrule\end{tabular}
\end{center}

\subsection{Two sources, read for what they actually say}

Item 20 is stated the way it is because Aspinwall's review is explicit about how far the
duality is currently understood: \emph{``This story began with the papers of [103, 96]
and, as we shall see, the full picture is still to be uncovered\ldots\ The exact number
[of topological classes of Calabi--Yau threefolds] is not known since a classification
remains elusive''} \source{Asp, ll.~3080--3087}. A Lean statement of the full duality
conjecture is therefore not available to formalize; the tractable first step is to state
and check \emph{one} numerical consequence the duality would imply (for instance, that the
Hodge numbers of a specific elliptically-fibered Calabi--Yau match a specific $K3\times
T^2$ moduli count) as a finite arithmetic check, in the spirit of
\cref{sec:atlas:examples}'s worked numeric examples, without claiming the general duality.

Item 21 is stated with the same care as item 2: GPR's own text says the semidirect-product
formula for $\Gamma$ ``holds true in the case of a single modulus'' and was only
\emph{proposed}, not derived in general, by the authors of a cited reference
\source{GPR, ll.~5432--5448}. Any Lean statement of item 21 for more than one modulus
would be formalizing a conjecture the source itself does not claim to have proved --
acceptable only if labeled \TC\ and attributed to the conjecturing paper, never presented
as \TL.

\subsection{Worked example with numbers: orbits of $24$ points}\label{sec:open:orbits}

Item 22 is the one entry in this section with an actual number to check by hand. EOT
observe that a finite group $G$ of symmetries of a K3 sigma model acts on $H^2(X,\Z)$
preserving at least $H^0$, $H^4$, $H^{2,0}$, $H^{0,2}$ and the K\"ahler form, so the
subspace of $H^{1,1}$ fixed by $G$ has dimension at least five, forcing the corresponding
action on $24$ points (via the Mathieu embedding) to split into \textbf{at least five
orbits} \source{EOT, ll.~last page, "$G$ cannot be $M_{24}$ itself"}. Their example is the
symmetry group of the Fermat quartic $X^4+Y^4+Z^4+W^4=0$ in $\mathbb{CP}^3$, the group
$(\Z_4)^2\rtimes S_4$ of order $384$, whose action on $24$ points splits into orbits of
lengths
\begin{equation}\label{eq:open:orbitcheck}
  1,\;1,\;2,\;4,\;16,\qquad 1+1+2+4+16=24,
\end{equation}
exactly five orbits, saturating the bound \source{EOT, final page}. This is a two-line
arithmetic check anyone can reproduce (\cref{eq:open:orbitcheck} is the whole content), and
it is exactly the kind of fact this project's tier-3 local prover could close today:
\begin{leancode}
-- not yet proved
theorem fermatQuartic_orbitLengths_sum :
    (1 : ℕ) + 1 + 2 + 4 + 16 = 24 := by decide
theorem fermatQuartic_five_orbits :
    ([1, 1, 2, 4, 16] : List ℕ).length = 5 := by decide
\end{leancode}
What is \emph{not} a $\star$ project is the theorem these two lines merely illustrate: that
\emph{every} subgroup $G\le M_{24}$ with at least five orbits on $24$ points arises this
way from an actual K3 symmetry, and conversely. That direction needs the global Torelli
theorem for K3 surfaces (cited but not formalized at \cref{ch:k3}) and an explicit
correspondence between $H^{1,1}$-fixed subspaces and $24$-point permutation
representations that this project has not built; it belongs at difficulty $\star\star\star$
and is the natural target for anyone wanting to extend \cref{ch:m24,ch:moonshine} in the
direction EOT themselves point.

%=====================================================================================
\section{After the first edition: Stream~3 and an external review}\label{sec:open:after}

This section was added after the chapters above were written. It records two things: a partial
answer to item 19 from a third companion library, and five recommendations from an external review of
this book and its papers, mapped onto the roadmap.

\subsection{Item 19 after Stream~3}\label{sec:open:stream3}
The library \texttt{DualScaleCosmology} (described in \texttt{docs/STREAM3\_WORKFLOW.md}) tests one
reading of the dual-scale proposal: a microscopic scale $\ell_{\rm micro}\sim\ell_P$ paired with a
macroscopic scale $\ell_{\rm macro}\sim c/H_0$. Its results, with their tiers:
{\sloppy\begin{itemize}
  \item Taking $\ell_P$ as the ultraviolet cutoff at the Hubble scale fails the Cohen--Kaplan--Nelson
    bound $L^3\Lambda^4\lesssim L M_P^2$ \source{\texttt{hep-th\_9803132.txt}, ll.~77--78} by more
    than thirty orders of magnitude: CKN's own horizon-scale cutoff is $\Lambda\sim10^{-2.5}$\,eV
    \source{\texttt{hep-th\_9803132.txt}, ll.~113--114}, against a Planck energy of about
    $1.2\times10^{28}$\,eV. \TA{} for the arithmetic
    (\lean{DualScaleCosmology.}\allowbreak\lean{CKNInstance.}\allowbreak\lean{planckEnergy\_gt\_}\allowbreak\lean{ckn\_horizon\_cutoff}), \TL{} for the
    two cited numbers.
  \item Reading $\ell_P$ and $c/H_0$ instead as a T-dual pair, with $\alpha'=\ell_P\cdot c/H_0$
    (\TC{}), the CKN bound holds exactly when the ultraviolet length is at least the self-dual length
    $s=\sqrt{\ell_P\,c/H_0}\approx47\,\mu$m
    (\lean{DualScaleCosmology.}\allowbreak\lean{SelfDualCutoff.}\allowbreak\lean{ckn\_iff\_uvLength\_}\allowbreak\lean{ge\_selfDual}, \TA{} as algebra).
    With $\rho_\Lambda=\Omega_\Lambda\cdot3H_0^2/(8\pi G)$ one has
    $\rho_\Lambda^{-1}=(8\pi/3\Omega_\Lambda)\,s^4$
    (\lean{DualScaleCosmology.}\allowbreak\lean{DarkEnergyScale.}\allowbreak\lean{rhoLambda\_inv\_eq}, \TA): the self-dual length
    \emph{is} the dark-energy length $\Lambda^{-1/4}\approx88\,\mu$m up to a fixed factor. Because it
    is one formula, the numerical agreement is algebra, not evidence.
  \item Read as the string's Regge slope, $\alpha'=s^2$ puts string excitations near $4$\,meV; the CMS
    dijet search excludes string resonances below $7.9$\,TeV in the models it tests
    \source{\texttt{1911\_03947.txt}, ll.~1044--1048}, a gap of more than $10^{30}$ in $\alpha'$
    (\lean{DualScaleCosmology.}\allowbreak\lean{CosmicString.}\allowbreak\lean{selfDual\_alphaPrime\_}\allowbreak\lean{exceeds\_cms\_ceiling}). Read as the
    radius of one extra dimension, $s$ lies above the torsion-balance bound of $30\,\mu$m
    \source{\texttt{2002\_11761.txt}, ll.~285--289} by a factor of about $1.6$ only, well within the
    order-one conventions involved: disfavored, not excluded. \TC{} for both readings.
\end{itemize}}
So item 19 now has a partial and mostly negative answer: the literal reading is excluded, and the
reading that survives adds nothing observable beyond existing proposals for a single micron-scale
extra dimension. What the machine checked here is algebra and arithmetic on cited constants; every
identification with physics is \TC.

\subsection{Recommendations from an external review}\label{sec:open:review}
An external review of this book and its companion papers made five recommendations for turning the
project into a convincing demonstration that theoretical physics can be machine-checked. Two coincide
with items above: replacing hard-coded Gram matrices by lattices defined as free $\Z$-modules with a
symmetric bilinear form is item~1 (with items~3--4 for the groups acting on them), and removing
theorems proved several times in different libraries is the atlas debt of items~12--16. The other
three are new.

\begin{center}\small
\begin{tabular}{@{}p{0.6cm}p{6.6cm}p{2.1cm}p{0.9cm}p{3.1cm}@{}}\toprule
\# & Item & Extends & Diff. & State of machinery \\ \midrule
23 & Make the tiers a compile-time property: attributes marking a declaration Tier A, L or C, and a
linter that walks the dependency graph and rejects a Tier A declaration that depends on a Tier L or C
one & \cref{ch:architecture,ch:pipeline} & $\star\star$ & Lean~4 attributes and environment
extensions exist; the dependency walk exists as \texttt{tools/}\allowbreak\texttt{lean\_depgraph.lean} \\
24 & An anti-vacuity linter: reject a theorem whose statement is \texttt{True}, a reflexive equation, or
does not mention one of its hypotheses' variables & \cref{ch:pipeline} & $\star$ & today done by hand:
the statement lock (\cref{ch:architecture}) and, in Stream~3, explicit counterexample theorems such as
\lean{DualScaleCosmology.}\allowbreak\lean{DualTower.}\allowbreak\lean{dualTower\_sum\_ge\_}\allowbreak\lean{needs\_product} \\
25 & A first continuous theory with no arithmetic shadow: the classical equivalence of the Nambu--Goto
and Polyakov actions, by eliminating the worldsheet metric through its equation of motion
\source{Tong §1.3, ll.~1255--1280} & \cref{ch:classical} & $\star\star\star$ & pointwise, a statement
about $2\times2$ matrices provable today; the full field-level statement needs variational calculus on
maps from a surface, which Mathlib does not yet provide \\
\bottomrule\end{tabular}
\end{center}

Item~25 has a finite-dimensional first step worth stating. At one point of the worldsheet let
$\gamma_{\alpha\beta}$ be the induced metric, with $\det\gamma<0$. The Polyakov Lagrangian density
$-\tfrac12\sqrt{-\det h}\,h^{\alpha\beta}\gamma_{\alpha\beta}$ is unchanged under $h\mapsto\lambda h$
(in two dimensions $\sqrt{-\det h}$ scales as $\lambda$ and $h^{\alpha\beta}$ as $\lambda^{-1}$), and at
$h=\gamma$ it equals $-\sqrt{-\det\gamma}$, the Nambu--Goto density. Both facts are identities about
$2\times2$ real matrices and could be kernel-checked with Mathlib as it stands; the calculus-of-variations
statement that this is where the Polyakov action is stationary is the part that is not yet available.
The Hubble-parameter theorem \lean{DualScaleCosmology.}\allowbreak\lean{ScaleFactorDuality.}\allowbreak\lean{hubble\_dual} shows that
pointwise derivative calculus of this kind is already routine in the project.

%=====================================================================================
\section{What the machine has checked}\label{sec:open:lean}

Every leancode block in this chapter marked \textbf{not yet proved} is exactly that: no
tier of the pipeline of \cref{ch:pipeline} has attempted it, and none of it is \TA. What
\emph{is} \TA, and what every item above explicitly builds from, are the two results
boxed below.

\begin{leanbox}[The base case every lattice item extends]
\begin{leancode}
theorem hyperbolicU_evenDiag : IsEvenDiag hyperbolicU := by ...
theorem hyperbolicU_unimodular : IsUnimodular hyperbolicU := ...
\end{leancode}
\textbf{Reading.} The $d=1$ Narain Gram matrix is even and unimodular.
\textbf{Proof idea.} Direct computation on an explicit $2\times2$ integer matrix.
\textbf{Not proved.} Anything for $d>1$; \cref{sec:open:narain2} is the honest gap between
this and item 1. \TA
\end{leanbox}

\begin{leanbox}[The base case every Mukai-generation item extends]\label{ex:open:e8weyl}
\begin{leancode}
theorem reflection_isometry (L : Gram n) (hL : Lᵀ = L) (v : Fin n → ℤ)
    (hv : latticeNorm L v = -2) :
    (reflection L v)ᵀ * L * reflection L v = L := by ...
theorem e8Neg_weyl_isometry (i : Fin 8) :
    (reflection e8Neg (Pi.single i 1))ᵀ * e8Neg * reflection e8Neg (Pi.single i 1)
      = e8Neg := ...
\end{leancode}
\textbf{Reading.} Each of the $8$ simple-root reflections is an isometry of $E_8(-1)$.
\textbf{Proof idea.} The general $(-2)$-vector reflection isometry, instantiated at each
of the $8$ standard basis vectors \source{Reflection.lean, ll.~128--181}.
\textbf{Not proved.} That the $8$ simple reflections \emph{generate} the full Weyl group of
$240$ roots -- stated verbatim in the file's own docstring as out of scope
\source{Reflection.lean, ll.~57--64} -- which is exactly item 4 above. \TA
\end{leanbox}

\begin{center}\small
\begin{tabular}{@{}p{6.3cm}cp{6.5cm}@{}}\toprule
Claim & Tier & Where \\ \midrule
The two boxed prerequisites, as stated & \TA & files cited; axioms \lean{propext},
  \lean{Classical.choice}, \lean{Quot.sound} only \\
Every leancode block marked "not yet proved" in this chapter & --- & not attempted by
  any tier of \cref{ch:pipeline}; not a claim of any kind \\
$A_n$ positivity and integrality for all $n$ & \TL & established in the literature since
  EOT's original conjecture (\cref{sec:open:moonshine-worked}), not re-derived here \\
GPR's $d=2$ duality group $\SLZ\times\SLZ\otimes_S[\Z_2\times\Z_2]$ in full & \TL & stated,
  not proved, in the source at the point read (\cref{sec:open:duality-group}) \\
Heterotic/type~II duality on $K3\times T^2$/Calabi--Yau (item 20) & \TC & Aspinwall's own
  assessment, "the full picture is still to be uncovered" \\
Six proofs of $\chi(K3)=24$ being "debt," and the remaining atlas clusters likewise &
  \TC & an engineering judgement, argued at \cref{ch:atlas} and \cref{sec:open:atlas} \\
\bottomrule\end{tabular}
\end{center}

\begin{summarybox}
\begin{itemize}[leftmargin=*,itemsep=1pt]
\item Twenty-two items, four groups: lattice/duality extensions (1--5), DFT/modular-form/
  moonshine targets (6--11), atlas engineering debt (12--16), a faithful non-shadow
  tadpole formalization (17--18), and physics-facing questions (19--22).
\item Difficulty is not the same as importance: item 8 ($A_n$ positivity past $n=2$) is
  $\star$ because the arithmetic is trivial once the $M_{24}$ character table entry is
  looked up correctly; item 3 (full $O(\Gamma^{4,20})$ action) is $\star\star\star$ because
  no group-generation machinery exists in the project at all.
\item Read the source before stating the target: items 2 and 21 both extend a Tier A
  result to a case the cited literature itself only conjectures or proves in a narrower
  case than this book's earlier chapters might suggest.
\item Some gaps are gaps in Mathlib (Siegel--Narain theta functions, Jacobi forms,
  singular cohomology of a compact 4-manifold), not gaps in effort; items 7, 10, 17--18
  name the missing library, not just the missing proof.
\item The atlas of \cref{ch:atlas} is itself a source for this chapter: five more
  unification-candidate clusters (items 12--16) are refactor projects of the exact shape
  worked out there in detail for three other clusters.
\end{itemize}
\end{summarybox}

\section*{Exercises}

\begin{exercise}[$\star$]
Carry out \cref{sec:open:narain2}'s block-diagonal argument formally: state and prove
\texttt{narain\_evenDiag\_two} and \texttt{narain\_unimodular\_two} in a scratch file
importing the module \texttt{Hyperbolic.lean} of \cref{ch:lattices}, using
\lean{Matrix.fromBlocks} and the matching block-determinant lemma (locate it yourself in
Mathlib). Report which tier of \cref{ch:pipeline} you think would have closed each goal.
\end{exercise}

\begin{exercise}[$\star$]
Verify \cref{eq:open:orbitcheck} by hand and state the two \lean{decide}-able lemmas of
\cref{sec:open:orbits} in a scratch file with no imports beyond core Lean. Then look up
one more subgroup of $M_{24}$ acting on $24$ points with at least five orbits (Mukai's or
Kondo's papers, cited at \cref{ch:m24}) and state the analogous two lemmas for it.
\end{exercise}

\begin{exercise}[$\star$]
Item 12 lists three Golay-code theorems as a unification candidate. Read the first pair
named there -- one in \texttt{Lean5Corpus}, one in \texttt{StringTheoryFoundation} -- state
in one sentence each what they actually prove, and decide whether they are the same fact
(as \cref{sec:atlas:unify}'s $\chi(K3)=24$ cluster turned out to be) or two different facts
that happen to both mention the number $24$.
\end{exercise}

\begin{exercise}[$\star\star$]
Read \cref{sec:open:duality-group}'s pinned excerpt of GPR's $d=2$ statement in the source
file itself (\texttt{papers/foundations/giveon\_hep-th\_9401139.txt}, around l.~7389).
State precisely, as a Lean proposition using the generators of \cref{ch:odd}, the single
piece this book already proves (the commuting subgroup), and explain in one paragraph what
additional generator and relation would be needed to reach the full $\Z_2\times\Z_2$
extension GPR states.
\end{exercise}

\begin{exercise}[$\star\star$]
Attempt item 6: state a Lean proposition capturing "$\mathcal D(G)=2d$ if and only if $G=1$
up to conjugation by an orthogonal matrix" as precisely as you can, using
\lean{dualScale\_ge}'s proof (\cref{ch:dualscale}) as a guide to where equality in the
trace-nonnegativity step forces $(G-1)G^{-1}(G-1)=0$. What extra hypothesis on $G$, if any,
does your statement need that \lean{dualScale\_ge} does not?
\end{exercise}

\begin{exercise}[$\star\star$]
Read \texttt{DualScaleStream2/Lattice/Reflection.lean}'s own "Not proved here" paragraph
(ll.~51--64) in full. List every one of its four explicitly named gaps, and for each say
whether it is a mathematics gap (a genuinely open or unformalized fact), an engineering
gap (missing Lean/Mathlib machinery), or an interpretation gap (a Lean object not yet
linked to the physical or geometric object it is meant to model).
\end{exercise}

\begin{exercise}[$\star\star$]
Choose one item from \cref{sec:open:atlas} (items 12--16) and carry out the atlas
refactor of \cref{sec:atlas:unify} for it: pick a canonical definition, state the other
theorem(s) in the cluster as corollaries, and re-run
\texttt{python3 tools/theorem\_atlas.py} to confirm the pair moves from the
unification-candidate list to the bridge list.
\end{exercise}

\begin{exercise}[$\star\star\star$]
Item 9 asks for $A_n$ as a virtual $M_{24}$ character. Using the character table excerpt
of Gaberdiel--Hohenegger--Volpato (\source{GHV §3.1, ll.~584--700}), derive by hand the
signed sum of irreducible characters giving $A_3=770$, verify it numerically, and state
(without necessarily proving) the corresponding Lean proposition in the style of
\cref{sec:open:moonshine-worked}'s \texttt{A3\_decomposition}, this time with the correct
coefficients.
\end{exercise}

\begin{exercise}[$\star\star\star$]
Pick any $\star\star\star$ item in this chapter, spend no more than two hours reading the
literature it depends on, and write a half-page assessment (not a proof) of the single
largest missing piece of machinery -- in Mathlib, in this project, or in the cited
literature -- standing between the current state and a Lean statement. Compare your
assessment to this chapter's stated "Mathlib / project state" column for that item.
\end{exercise}

\section*{Further reading}

The improvements and next targets this chapter draws from are \texttt{LL.md §S2-I}
(items 1--6, process; item 7, the four Stream~3 mathematics targets behind items 1, 3, 7
and 9 above) and the atlas's own unification-candidate list,
\texttt{papers/book/generated/atlas.md} (behind items 12--16, and behind
\cref{ch:atlas}'s three worked clusters this chapter deliberately does not repeat). The
three literature sources read directly for this chapter's physics-facing section are
GPR \source{giveon\_hep-th\_9401139.txt, ll.~5420--5450, 7380--7395, 2905--2925}, Aspinwall
\source{aspinwall\_hep-th\_9611137.txt, ll.~3075--3090}, and EOT
\source{eguchi\_ooguri\_tachikawa\_1004\_0956.txt, ll.~175--215 and the final page}; each
excerpt was read in the source file itself, not from a secondary summary, per this book's
own rule (\S1 of the Book Bible: read sources, do not write from memory). For how any of
these items would actually be attempted --
statement design first, symbolic pre-verification, tiered proof search, and a gate that
recompiles every claim -- see \cref{ch:pipeline}; for what "done" measurably looks like for
an atlas item specifically, see \cref{ch:atlas}'s closing remark that re-running the
atlas script should shrink the unification list, one refactored cluster at a time.
